% Part: sets-functions-relations
% Chapter: sets
% Section: important-sets

\documentclass[../../../include/open-logic-section]{subfiles}

\begin{document}

\olfileid{sfr}{set}{imp}
\olsection{కొన్ని ముఖ్యమైన సమితులు}

\begin{ex}
మన పరిశీలనలో ఎక్కువగా గణిత వస్తువులు \tetoken{మూలకాలుగా}{!!{element}s} గల సమితులే
ఉంటాయి. వాటిలో నాలుగు సమితులు ప్రత్యేక పేర్లతో పిలిచేంత ముఖ్యమైనవి:
\begin{multline*}
    \Nat = \{0, 1, 2, 3, \ldots\} \\
    \shoveright{\text{సహజ సంఖ్యల సమితి}}\\
    \shoveleft{\Int = \{\ldots, -2, -1, 0, 1, 2, \ldots\}} \\
    \shoveright{\text{పూర్ణ సంఖ్యల సమితి}}\\
    \shoveleft{\Rat = \Setabs{\nicefrac{m}{n}}{m, n \in \Int\text{ మరియు }n \neq 0}}\\
    \shoveright{\text{కరణీయ సంఖ్యల సమితి}}\\
    \shoveleft{\Real = (-\infty, \infty)}\\
    \text{వాస్తవ సంఖ్యల సమితి (అవిచ్ఛిన్న సమితి)}
\end{multline*}
ఇవన్నీ \emph{అనంత} సమితులు; అంటే ప్రతి సమితిలోనూ
అనంతంగా \tetoken{మూలకాలు}{!!{element}s} ఉంటాయి.

ఈ సమితులను వరుసగా పరిశీలిస్తే, ఒక్కో దశలో మన సంఖ్యల సమూహానికి
\emph{మరిన్ని} సంఖ్యలను చేర్చుతున్నాం. నిజానికి
$\Nat \subseteq \Int \subseteq \Rat \subseteq \Real$ అని స్పష్టమే:
ప్రతి సహజ సంఖ్య పూర్ణ సంఖ్య; ప్రతి పూర్ణ సంఖ్య కరణీయ సంఖ్య;
ప్రతి కరణీయ సంఖ్య వాస్తవ సంఖ్య. అదే విధంగా
$\Nat \subsetneq \Int \subsetneq \Rat$ అనేదీ స్పష్టమే.
ఎందుకంటే $-1$ పూర్ణ సంఖ్య అయినా సహజ సంఖ్య కాదు;
$\nicefrac{1}{2}$ కరణీయ సంఖ్య అయినా పూర్ణ సంఖ్య కాదు.
కానీ $\Rat \subsetneq \Real$ అనేది, అంటే కరణీయాలు కాని
వాస్తవ సంఖ్యలు ఉన్నాయనేది, అంత సులభంగా స్పష్టం
కాదు\oliflabeldef{sfr:arith:real:realline}{; ఈ విషయాన్ని
\olref[arith][real]{realline}లో తిరిగి పరిశీలిస్తాం}{}.

కొన్నిసార్లు ధన పూర్ణ సంఖ్యల సమితి $\PosInt = \{1, 2, 3, \dots\}$ను,
మొదటి రెండు సహజ సంఖ్యలు మాత్రమే గల $\Bin = \{0, 1\}$ సమితిని కూడా వాడతాం.
\end{ex}

\begin{tagblock}{compsci}
\begin{ex}[సంకేతమాలలు]
మరో ఆసక్తికరమైన ఉదాహరణ, $A$ అనే వర్ణమాలపై గల
\emph{పరిమిత సంకేతమాలల} సమితి $A^{*}$. $A$లోని మూలకాలతో
ఏర్పడే ఏ పరిమిత క్రమాన్నైనా $A$పై ఒక సంకేతమాల అంటాం.
ప్రతి వర్ణమాల $A$కూ, $A$పై గల సంకేతమాలలలో
\emph{శూన్య సంకేతమాల $\Lambda$}ను కూడా చేర్చుతాం. ఉదాహరణకు,
\begin{multline*}
\Bin^*
=\{\Lambda,0,1,00,01,10,11,\\
000,001,010,011,100,101,110,111,0000,\ldots\}.
\end{multline*}
$x=x_{1}\ldots x_{n}\in A^{*}$ అనేది $A$లోని $n$ ``అక్షరాలతో''
ఏర్పడిన సంకేతమాల అయితే, దాని \emph{పొడవు} $n$ అని చెప్పి
$\len{x}=n$ అని రాస్తాం.
\end{ex}
\end{tagblock}

\begin{ex}[అనంత క్రమాలు]
ఏ సమితి $A$కైనా, $A$లోని \tetoken{మూలకాల}{!!{element}s} అనంత క్రమాల
సమితి $A^\omega$ను కూడా పరిగణించవచ్చు. $a_1a_2a_3a_4\dots$
అనే అనంత క్రమం, ఒక దిశలో అనంతంగా కొనసాగే వస్తువుల జాబితా;
అందులో ప్రతి వస్తువూ $A$లో \tetoken{ఒక మూలకం}{!!a{element}}.
\end{ex}

\end{document}
